\chapter{LU and PLU Factorisation}

Gaussian elimination is a sequence of row operations. Chapter~3 showed that each row operation is a multiplication by an elementary matrix, so elimination is a matrix factorisation in disguise. Recording the factorisation, instead of discarding the multipliers, turns elimination into a reusable tool: factor once, then solve with many right-hand sides cheaply. This chapter develops the $LU$ and $PLU$ factorisations, the Cholesky factorisation of symmetric positive definite matrices, and the cost model that explains why numerical linear algebra is organised around factorisations.

\section{Elimination as a product of matrices}

Take an $n\times n$ matrix $A$ and suppose elimination to upper triangular form needs no row swaps. At step $j$ we subtract multiples of row $j$ from the rows below it. The multiple used for row $i>j$ is
\[ \ell_{ij}=\frac{a^{(j)}_{ij}}{a^{(j)}_{jj}}, \]
where $a^{(j)}_{ij}$ is the current entry. The pivot $a^{(j)}_{jj}$ must be nonzero.

\begin{example}
\[
A=\begin{pmatrix}2&1&1\\4&3&3\\8&7&9\end{pmatrix}.
\]
Step 1: $\ell_{21}=2$ and $\ell_{31}=4$, giving
$\begin{pmatrix}2&1&1\\0&1&1\\0&3&5\end{pmatrix}$. Step 2: $\ell_{32}=3$, giving the upper triangular matrix
\[ U=\begin{pmatrix}2&1&1\\0&1&1\\0&0&2\end{pmatrix}. \]
The multipliers are $\ell_{21}=2,\ \ell_{31}=4,\ \ell_{32}=3$. Remarkably,
\[ A=\begin{pmatrix}1&0&0\\2&1&0\\4&3&1\end{pmatrix}\begin{pmatrix}2&1&1\\0&1&1\\0&0&2\end{pmatrix}=LU, \]
where $L$ is the lower triangular matrix with the multipliers placed below a diagonal of ones.
\end{example}

\begin{theorem}[$LU$ factorisation]\label{thm:lu}
Suppose Gaussian elimination of $A\in\R^{n\times n}$ to upper triangular form $U$ succeeds without row exchanges. Then $A=LU$, where $U$ is upper triangular with the pivots on its diagonal and $L$ is unit lower triangular (ones on the diagonal) whose $(i,j)$ entry for $i>j$ is the multiplier $\ell_{ij}$.
\end{theorem}

\begin{proof}
Step $j$ of elimination is multiplication on the left by $E_j=I-\boldsymbol{\ell}_j\mathbf{e}_j\T$, where $\boldsymbol{\ell}_j=(0,\dots,0,\ell_{j+1,j},\dots,\ell_{nj})\T$. Thus $E_{n-1}\cdots E_1A=U$ and $A=E_1^{-1}\cdots E_{n-1}^{-1}U$. Each inverse is $E_j^{-1}=I+\boldsymbol{\ell}_j\mathbf{e}_j\T$ (check: $(I-\boldsymbol\ell\mathbf{e}\T)(I+\boldsymbol\ell\mathbf{e}\T)=I-\boldsymbol\ell(\mathbf{e}\T\boldsymbol\ell)\mathbf{e}\T=I$ since $\mathbf{e}_j\T\boldsymbol\ell_j=0$). Multiplying the inverses in the order $E_1^{-1}E_2^{-1}\cdots$, a cross term such as $\boldsymbol\ell_j\mathbf{e}_j\T\boldsymbol\ell_k\mathbf{e}_k\T$ with $j<k$ vanishes, because $\boldsymbol\ell_k$ has zeros in positions $1,\dots,k$ and so $\mathbf{e}_j\T\boldsymbol\ell_k=0$. Hence the product is $I+\sum_j\boldsymbol\ell_j\mathbf{e}_j\T=L$. This is the key point: unlike $E_{n-1}\cdots E_1$, the product of the inverses in this order simply \emph{collects} the multipliers with no interaction.
\end{proof}

\subsection*{Solving $A\mathbf{x}=\mathbf{b}$ with $LU$}

Given $A=LU$, the system $A\mathbf{x}=\mathbf{b}$ becomes two triangular systems:
\[ L\mathbf{y}=\mathbf{b}\quad(\text{forward substitution}),\qquad U\mathbf{x}=\mathbf{y}\quad(\text{back substitution}). \]
Each triangular solve costs about $n^2$ operations, against about $\tfrac23n^3$ for the factorisation. Once $A$ is factored, every further right-hand side costs only $O(n^2)$.

\begin{example}
Solve $A\mathbf{x}=(4,10,22)\T$ for the matrix above. Forward: $y_1=4$, $2y_1+y_2=10\Rightarrow y_2=2$, $4y_1+3y_2+y_3=22\Rightarrow y_3=0$. Back: $2x_3=0\Rightarrow x_3=0$; $x_2+x_3=2\Rightarrow x_2=2$; $2x_1+x_2+x_3=4\Rightarrow x_1=1$. So $\mathbf{x}=(1,2,0)\T$.
\end{example}

\begin{theorem}[Uniqueness of $LU$]\label{thm:luuniq}
If $A=L_1U_1=L_2U_2$ with $L_i$ unit lower triangular and $U_i$ upper triangular and invertible, then $L_1=L_2$ and $U_1=U_2$.
\end{theorem}

\begin{proof}
$L_2^{-1}L_1=U_2U_1^{-1}$. The left side is unit lower triangular (products and inverses of such matrices are again unit lower triangular), the right side is upper triangular. A matrix that is both is diagonal, and with unit diagonal it must equal $I$.
\end{proof}

\section{When pivots vanish: $PA=LU$}

Elimination may meet a zero pivot. For $A=\begin{pmatrix}0&1\\1&1\end{pmatrix}$ the first pivot is $0$, yet $A$ is invertible; swapping the rows repairs this. A tiny pivot is nearly as bad: dividing by it produces huge multipliers and amplifies rounding error.

\begin{definition}[Partial pivoting]
At step $j$, choose as pivot the entry of largest absolute value among $a^{(j)}_{jj},\dots,a^{(j)}_{nj}$ and swap its row into position $j$.
\end{definition}

Partial pivoting keeps every multiplier $\abs{\ell_{ij}}\le1$.

\begin{theorem}[$PA=LU$]\label{thm:plu}
Every invertible matrix $A$ admits a factorisation $PA=LU$ with $P$ a permutation matrix, $L$ unit lower triangular with entries of absolute value at most $1$, and $U$ upper triangular and invertible. (Every square matrix admits such a factorisation with $U$ not necessarily invertible.)
\end{theorem}

\begin{proof}[Sketch]
Perform elimination with row swaps. The swaps can be moved to the front: performing all swaps first and then eliminating gives the same multipliers up to reordering, because swapping two rows of the matrix being reduced swaps the corresponding rows of the multiplier matrix accumulated so far. Collect the swaps into $P$. The choice of the largest pivot makes each multiplier at most $1$ in absolute value. If $A$ is invertible, a nonzero pivot exists at every step, since otherwise the whole column below and including the diagonal would vanish, contradicting invertibility.
\end{proof}

To solve $A\mathbf{x}=\mathbf{b}$: solve $L\mathbf{y}=P\mathbf{b}$, then $U\mathbf{x}=\mathbf{y}$.

\begin{nuance}{why pivoting matters numerically}
Consider $\begin{pmatrix}10^{-20}&1\\1&1\end{pmatrix}\mathbf{x}=\begin{pmatrix}1\\2\end{pmatrix}$. Without pivoting, the multiplier is $10^{20}$; the second equation becomes $(1-10^{20})x_2=2-10^{20}$, and in double precision both coefficients round to $-10^{20}$, giving $x_2=1$ and then $x_1=0$, whereas the true solution is $x_1\approx x_2\approx1$. With a row swap the multiplier is $10^{-20}$ and the answer is accurate. In practice, partial pivoting is stable for nearly all matrices; examples of exponential growth exist but are extremely rare.
\end{nuance}

\section{The cost of factorisations}

Counting a multiply-add as one operation:
\begin{center}
\begin{tabular}{ll}
\toprule
Task & Operations \\
\midrule
Compute $LU$ or $PA=LU$ & $\tfrac13n^3$ \\
Forward or back substitution & $\tfrac12n^2$ each \\
Solve $A\mathbf{x}=\mathbf{b}$ for $k$ right-hand sides & $\tfrac13n^3+kn^2$ \\
Compute $A^{-1}$ & $n^3$ \\
Multiply two $n\times n$ matrices & $n^3$ \\
\bottomrule
\end{tabular}
\end{center}

The determinant is a by-product: $\det A=\pm\prod_i u_{ii}$, the sign being that of the permutation (Chapter~5). This is how software actually computes determinants, not by cofactor expansion, which costs $n!$ operations.

\begin{interview}{factor once, solve many times}
Typical question: ``You must solve $A\mathbf{x}_k=\mathbf{b}_k$ for a thousand different $\mathbf{b}_k$ with the same $A$. How?'' Compute the $PA=LU$ factorisation once, $O(n^3)$, then each solve costs $O(n^2)$. Computing $A^{-1}$ and multiplying is both slower and less accurate. The same principle underlies implicit time-stepping of differential equations and Newton's method with a fixed Jacobian.
\end{interview}

\section{Symmetric positive definite matrices: Cholesky}

Chapter~16 studies positive definite matrices in depth; here we need only the definition. A symmetric matrix $A$ is \emph{positive definite} if $\mathbf{x}\T A\mathbf{x}>0$ for every $\mathbf{x}\ne\mathbf{0}$.

\begin{theorem}[Cholesky factorisation]\label{thm:chol}
A symmetric matrix $A$ is positive definite if and only if $A=GG\T$ for some lower triangular matrix $G$ with positive diagonal entries. The factor $G$ is unique.
\end{theorem}

\begin{proof}
If $A=GG\T$ with $G$ invertible then $\mathbf{x}\T A\mathbf{x}=\norm{G\T\mathbf{x}}^2>0$ for $\mathbf{x}\ne\mathbf{0}$. Conversely, suppose $A$ is positive definite. Each leading principal submatrix is positive definite, hence invertible, so elimination without pivoting runs and gives $A=LU$ with positive pivots (the $k$-th pivot is the ratio of leading principal minors, which are positive). Because $A$ is symmetric, $L^{-1}AL^{-\top}=UL^{-\top}$ is both symmetric and upper triangular, hence diagonal: $U=DL\T$ with $D=\diag(u_{11},\dots,u_{nn})$ and all $u_{ii}>0$. So $A=LDL\T=(LD^{1/2})(LD^{1/2})\T$; set $G=LD^{1/2}$. Uniqueness follows from uniqueness of $LU$ (Theorem~\ref{thm:luuniq}).
\end{proof}

Cholesky uses half the work of $LU$ (about $\tfrac16n^3$), needs no pivoting, and is the standard way to solve normal equations and to sample from a multivariate Gaussian: if $\mathbf{z}\sim N(\mathbf{0},I)$ and $\Sigma=GG\T$, then $G\mathbf{z}\sim N(\mathbf{0},\Sigma)$.

\begin{example}
For $A=\begin{pmatrix}4&2\\2&3\end{pmatrix}$ we get $g_{11}=2$, $g_{21}=2/2=1$, $g_{22}=\sqrt{3-1}=\sqrt2$:
\[ G=\begin{pmatrix}2&0\\1&\sqrt2\end{pmatrix},\qquad GG\T=\begin{pmatrix}4&2\\2&3\end{pmatrix}. \]
\end{example}

\section{Other factorisations at a glance}

The factorisations of this book form a small family, each exposing different structure.

\begin{center}
\begin{tabular}{lll}
\toprule
Factorisation & Form & Chapter \\
\midrule
LU / PLU & $PA=LU$ & 4 \\
Cholesky & $A=GG\T$ & 4, 16 \\
QR & $A=QR$ & 11 \\
Eigen / spectral & $A=V\Lambda V^{-1}$, $A=Q\Lambda Q\T$ & 13--15 \\
SVD & $A=U\Sigma V\T$ & 17 \\
\bottomrule
\end{tabular}
\end{center}

\begin{summarybox}
\begin{itemize}
  \item Elimination is $E_{n-1}\cdots E_1A=U$; the multipliers collect into a unit lower triangular $L$ with $A=LU$, unique when it exists.
  \item Solving with $LU$ is two triangular solves of cost $O(n^2)$ after an $O(n^3)$ factorisation; factor once, solve many times.
  \item Row swaps give $PA=LU$; partial pivoting keeps multipliers at most $1$ and protects against rounding error.
  \item Symmetric positive definite matrices factor as $GG\T$ (Cholesky) at half the cost, with no pivoting.
  \item Never compute $A^{-1}$ or a cofactor determinant to solve a problem; use a factorisation.
\end{itemize}
\end{summarybox}

\section{Exercises}

\begin{exercise}
Compute the $LU$ factorisation of $\begin{pmatrix}1&2&1\\3&8&2\\2&4&5\end{pmatrix}$ and use it to solve $A\mathbf{x}=(1,0,2)\T$.
\end{exercise}

\begin{exercise}
Find $P$, $L$, $U$ with $PA=LU$ for $A=\begin{pmatrix}0&2&1\\1&1&0\\2&0&3\end{pmatrix}$.
\end{exercise}

\begin{exercise}
Show that the product of two unit lower triangular matrices is unit lower triangular, and that the inverse of one is unit lower triangular.
\end{exercise}

\begin{exercise}
Prove that a matrix with $a_{11}=0$ and $a_{21}\ne0$ cannot have an $LU$ factorisation with $L$ unit lower triangular and $U$ invertible, even if $A$ is invertible, and explain how this relates to $PA=LU$.
\end{exercise}

\begin{exercise}
Compute the Cholesky factor of $\begin{pmatrix}9&6&3\\6&5&4\\3&4&11\end{pmatrix}$.
\end{exercise}

\begin{exercise}
Show that if $A=LU$ and $\det A\ne0$ then $\det A=\prod u_{ii}$ (use that $\det$ of a triangular matrix is the product of its diagonal; this is proved in Chapter~5).
\end{exercise}

\begin{exercise}
The \emph{band} matrix with $a_{ii}=2$, $a_{i,i+1}=a_{i+1,i}=-1$ arises from discretising a second derivative. Find its $LU$ factors for $n=4$ and observe that $L$ and $U$ are also banded. Why does this make solving such systems cost only $O(n)$?
\end{exercise}

\begin{exercise}
Suppose $\Sigma=\begin{pmatrix}1&0.8\\0.8&1\end{pmatrix}$. Find $G$ with $\Sigma=GG\T$ and describe how to generate correlated Gaussian samples from independent ones.
\end{exercise}
